\begin{figure}[htbp]
    \centering
    \begin{tikzpicture}[scale=1.2, >=stealth]
        % Tô màu miền phương án \Omega trước để nằm dưới các nét vẽ
        \fill[cyan!15] (0,2) -- (2,4) -- (4,2) -- (3,0) -- (0,0) -- cycle;

        % Vẽ hệ trục tọa độ
        \draw[->, thick] (-1,0) -- (6,0) node[right] {$x_1$};
        \draw[->, thick] (0,-1) -- (0,5) node[above] {$x_2$};
        \node[below left] at (0,0) {$O$};

        % Vẽ các đường thẳng ràng buộc (viền của tập phương án)
        \draw[thick, blue] (-1,1) -- (5,7) node[right] {$(d_1)$}; 
        \draw[thick, red] (1,5) -- (5,1) node[right] {$(d_2)$};
        \draw[thick, green!60!black] (4.5,-1.5) -- (2.5,4.5) node[above] {$(d_3)$};

        % Đánh dấu các đỉnh của đa giác \Omega
        \filldraw (0,2) circle (1.5pt) node[left] {$A$};
        \filldraw (2,4) circle (1.5pt) node[above] {$B$};
        \filldraw (4,2) circle (1.5pt) node[right] {$C$};
        \filldraw (3,0) circle (1.5pt) node[below] {$D$};
        \node at (2,1.5) {\Large $\Omega$};

        % Vẽ đường mức \Delta và vector pháp tuyến n
        \draw[dashed, thick, orange] (-1,3) -- (5,-0.5) node[below] {$\Delta$ (Đường mức)};
        \draw[->, very thick, orange] (2,1.25) -- (2.8, 2.6) node[above right] {$\vec{n}$};

        % Vẽ đường mức tới hạn tại đỉnh B (PATƯ)
        \draw[dashed, thick, orange!70!black] (0.5,4.8) -- (5,2.2) node[right] {$\Delta_{max}$};
        \filldraw[orange!70!black] (2,4) circle (2pt);
    \end{tikzpicture}
    \caption{Minh họa tập phương án $\Omega$ và sự tịnh tiến đường mức $\Delta$ tìm đỉnh tối ưu}
    \label{fig:phuongphap_hinhhoc}
\end{figure}